% Part: first-order-logic
% Chapter: syntax-and-semantics
% Section: introduction

\documentclass[../../../include/open-logic-section]{subfiles}

\begin{document}

\olfileid{fol}{syn}{int}

\olsection{مقدمة}

لكي نطور نظرية منطق الرتبة الأولى ونظريته الفوقية، يجب أولًا أن نعرّف
تركيب تعبيراته ودلالتها. وتعبيرات منطق الرتبة الأولى حدود و!!{formula}s.
تُبنى الحدود من~!!{variable}s و!!{constant}s و!!{function}s. أما
!!^{formula}s فتُبنى من~!!{predicate}s مع الحدود، وهذه تكوّن أبسط
!!{formula}s، أي الصيغ «الذرية»، ثم يمكننا أن نبني من
!!{formula}s الذرية صيغًا أعقد باستخدام الروابط المنطقية والأسوار.
وتوجد طرائق كثيرة مختلفة لصياغة قواعد التكوين؛ ولا نقدم إلا طريقة
واحدة ممكنة. فقد تختار أنظمة أخرى رموزًا مختلفة، وتختار مجموعات مختلفة
من الروابط بوصفها أولية، وتستخدم الأقواس على نحو مختلف، أو لا تستخدمها
إطلاقًا كما في الترميز البولندي المسمى. لكن القاسم المشترك بين جميع
المقاربات هو أن قواعد التكوين تعرّف مجموعة الحدود و!!{formula}s
تعريفًا 
\emph{استقرائيًا}. وإذا أُحكم هذا التعريف، أمكن توليد كل تعبير، في
الجوهر، بطريقة واحدة فقط وفق قواعد التكوين. ولما كان التعريف
الاستقرائي ينتج تعبيرات~\emph{وحيدة القراءة}، أمكننا إسناد المعاني
إليها بالطريقة نفسها، أي بالتعريف الاستقرائي.

إسناد المعنى إلى التعبيرات هو ميدان الدلالة. والمفهوم المركزي في
الدلالة هو الإشباع في~!!a{structure}. ويُسند~!!^a{structure} معنى إلى
لبنات اللغة: فـ!!a{domain} مجموعة غير خالية من الأشياء. وتُفسر الأسوار
على أنها تتراوح على هذا المجال، وتُسند إلى~!!{constant}s عناصر من
المجال، ويُسند إلى كل رمز من~!!{function}s ذي رتبة~$n$ دالة من
$D^n$ إلى~$D$، حيث~$D$ هو~!!{domain}، وتُسند إلى~!!{predicate}s علاقات
على~!!{domain}. ويكوّن~!!{domain}، مع الإسنادات إلى المفردات الأساسية،
!!a{structure}. وقد تظهر~!!^{variable}s في~!!{formula}s، ولكي نقدم
دلالة لها يجب أيضًا أن نسند إليها~!!{element}s من~!!{domain}، وهذا هو
إسناد المتغيرات. وأخيرًا تجمع علاقة الإشباع هذه الأمور. قد تُشبَع
!!^a{formula} في~!!a{structure}~$\Struct{M}$ نسبة إلى إسناد متغيرات~$s$،
ونكتب ذلك~$\Sat{M}{!A}[s]$. وتُعرّف هذه العلاقة أيضًا بالاستقراء على
بنية~$!A$، باستخدام جداول الصدق للروابط المنطقية كي نعرّف، مثلًا،
إشباع~$!A \land !B$ بدلالة إشباع~$!A$ و$!B$ أو عدم إشباعهما. ويتبين
بعدئذ أن إسناد المتغيرات غير مهم إذا كانت~!!{formula}~$!A$
!!a{sentence}، أي لا متغيرات حرة فيها، ولذلك نستطيع الحديث عن إشباع
!!{sentence}s أو عدم إشباعها ببساطة في~!!{structure}s.

وبناء على علاقة الإشباع~$\Sat{M}{!A}$ للجمل، يمكننا تعريف المفاهيم
الدلالية الأساسية: الصحة واللزوم وقابلية الإشباع. تكون الجملة صحيحة،
$\Entails !A$، إذا أشبعتها كل بنية. وتلزم من مجموعة~!!{sentence}s،
$\Gamma \Entails !A$، إذا كان كل~!!{structure} يشبع جميع
!!{sentence}s في~$\Gamma$ يشبع~$!A$ أيضًا. وتكون مجموعة من الجمل قابلة
للإشباع إذا كان ثمة~!!{structure} يشبع كل~!!{sentence}s فيها في آن
واحد. ولأن~!!{formula}s معرّفة استقرائيًا، والإشباع معرّف بدوره
بالاستقراء على بنية~!!{formula}s، يمكننا استخدام الاستقراء لإثبات خواص
دلالتنا وربط المفاهيم الدلالية المعرّفة.




\end{document}
